% ============================================================
% CHAPTER 16: RANDOMIZED ROUNDING (matches Vazirani Ch. 16)
% ============================================================
\setcounter{chapter}{15}
\chapter{Randomized Rounding}
\label{ch:randomized_rounding}

\section{Intuition: LP Probabilistic Rounding}

\begin{intuitionbox}
\textbf{Peeling and Layering Intuition:} LP rounding algorithms are often deterministic (e.g. Set Cover rounding at threshold $1/f$). However, for problems like Max-SAT, deterministic thresholding can be suboptimal. Instead, we can interpret the fractional LP variables $y_i^*$ as \emph{probabilities}. By setting variable $x_i$ to True with probability $y_i^*$, we ensure that the expected value of the assignment is close to the LP optimal value. By combining this with a simple coin-flip assignment (which performs well on large clauses, while LP rounding performs well on small clauses), we obtain a $3/4$-approximation.
\end{intuitionbox}

\section{Problem Definition}
Given $n$ boolean variables $x_1, \dots, x_n$ and $m$ clauses $C_1, \dots, C_m$ with weights $w_j \ge 0$, the \emph{Maximum Satisfiability} (Max-SAT) problem is to find a truth assignment to the variables that maximizes the total weight of satisfied clauses.

The decision version of SAT is NP-complete. Max-SAT is NP-hard. We seek a randomized $(1 - 1/e)$-approximation (approx. 0.632) via LP rounding, and a $3/4$-approximation (0.75) by choosing the best of LP rounding and random coin-flipping.

\section{Algorithm and Python Implementation}
We formulate the Max-SAT problem as an integer program, solve its LP relaxation using our Simplex solver, and perform randomized rounding.

\begin{lstlisting}[style=python, caption=Randomized Rounding for Max-SAT]
import random


class Simplex:
    """Maximize c x  s.t. Ax <= b, x >= 0."""
    def __init__(self, A, b, c):
        self.m = len(A)
        self.n = len(A[0]) if A else 0
        self.A = [row[:] for row in A]
        self.b = b[:]
        self.c = c[:]

    def solve(self):
        tab = [self.A[i] + [1 if j == i else 0 for j in range(self.m)]
               + [self.b[i]] for i in range(self.m)]
        tab.append([-v for v in self.c] + [0] * self.m + [0])
        while True:
            e = min(range(self.n + self.m), key=lambda j: tab[-1][j])
            if tab[-1][e] >= -1e-9:
                break
            l = min(range(self.m), key=lambda i: tab[i][-1] / tab[i][e]
                    if tab[i][e] > 1e-9 else float('inf'))
            if tab[l][e] <= 1e-9:
                return None, float('inf')
            p = tab[l][e]
            tab[l] = [v / p for v in tab[l]]
            for i in range(self.m + 1):
                if i != l:
                    f = tab[i][e]
                    tab[i] = [tab[i][j] - f * tab[l][j]
                              for j in range(len(tab[0]))]
        x = [0.0] * self.n
        for i in range(self.m):
            for j in range(self.n + self.m):
                if (abs(tab[i][j] - 1) < 1e-6 and
                        sum(abs(tab[r][j]) for r in range(self.m + 1)) < 1.1):
                    if j < self.n:
                        x[j] = max(0.0, tab[i][-1])
        return x, tab[-1][-1]


def evaluate_assignment(assignment, clauses, weights):
    """Total weight of clauses satisfied by a boolean assignment."""
    total = 0.0
    for j, (pos, neg) in enumerate(clauses):
        satisfied = any(assignment[i] for i in pos)
        if not satisfied:
            satisfied = any(not assignment[i] for i in neg)
        if satisfied:
            total += weights[j]
    return total


def solve_max_sat_lp(n_vars, clauses, weights):
    """Formulate and solve LP relaxation for Max-SAT."""
    n_cols = n_vars + len(clauses)
    A, b = [], []
    
    # Clause constraints: z_j - sum_{i in C_j^+} y_i + sum_{i in C_j^-} y_i <= |C_j^-|
    for j, (pos, neg) in enumerate(clauses):
        row = [0.0] * n_cols
        row[n_vars + j] = 1.0
        for i in pos: row[i] = -1.0
        for i in neg: row[i] = 1.0
        A.append(row)
        b.append(float(len(neg)))
        
    # Bounds: y_i <= 1, z_j <= 1
    for i in range(n_vars):
        row = [0.0] * n_cols; row[i] = 1.0; A.append(row); b.append(1.0)
    for j in range(len(clauses)):
        row = [0.0] * n_cols; row[n_vars + j] = 1.0; A.append(row); b.append(1.0)
        
    c = [0.0] * n_cols
    for j in range(len(clauses)): c[n_vars + j] = weights[j]
    
    solver = Simplex(A, b, c)
    sol, obj = solver.solve()
    return sol[:n_vars], sol[n_vars:], obj

def randomized_rounding_max_sat(n_vars, clauses, weights, y_lp, trials=200):
    """Applies randomized rounding to LP variable values."""
    best_weight, best_assignment, total_weight = -1.0, [], 0.0
    for _ in range(trials):
        assignment = [random.random() < y_lp[i] for i in range(n_vars)]
        weight = evaluate_assignment(assignment, clauses, weights)
        total_weight += weight
        if weight > best_weight:
            best_weight, best_assignment = weight, assignment
    return best_assignment, total_weight / trials, best_weight
\end{lstlisting}

\section{Analysis: Approximations and Probability}
For a clause $C_j$ with size $k$, let its index set of positive and negative literals be $C_j^+$ and $C_j^-$. The LP relaxation includes:
\[
z_j \le \sum_{i \in C_j^+} y_i + \sum_{i \in C_j^-} (1 - y_i)
\]
Under randomized rounding, variable $x_i$ is set to True with probability $y_i^*$. The probability that clause $C_j$ is satisfied is:
\[
1 - \prod_{i \in C_j^+} (1 - y_i^*) \prod_{i \in C_j^-} y_i^*
\]
Using the arithmetic-geometric mean inequality, we can show this probability is at least:
\[
1 - \left(1 - \frac{z_j}{k}\right)^k
\]
Since $g(z) = 1 - (1 - z/k)^k$ is concave and $g(0) = 0, g(1) = 1 - (1 - 1/k)^k$, we have:
\[
\text{Pr}[C_j \text{ is satisfied}] \ge \left[1 - \left(1 - \frac{1}{k}\right)^k\right] z_j \ge \left(1 - \frac{1}{e}\right) z_j
\]
Summing over all clauses, the expected weight satisfied is $\ge (1 - 1/e) \cdot \text{OPT}_{LP} \approx 0.632 \cdot \text{OPT}$.
For small clauses, this bound is much higher (e.g. $k=1$ gives $1 \cdot z_j$, $k=2$ gives $0.75 \cdot z_j$). For large clauses, the coin-flip baseline satisfies them with probability $1 - 2^{-k}$. Choosing the best of both algorithms yields a $3/4$-approximation (Goemans-Williamson 1994).

\section{Applications}
\begin{itemize}
    \item \textbf{Satisfiability Solvers:} Finding satisfying assignments for weighted CNF formulas in hardware design verification.
    \item \textbf{Database Query Optimizers:} Selecting indexing configurations to satisfy as many query filters as possible under physical memory limits.
    \item \textbf{Digital Circuit Testing:} Generating test vectors that cover maximum gate combinations.
\end{itemize}

\subsection*{Real Example: Database Query Filter Maximization}
\textbf{Scenario:} A high-speed database system (e.g., Google Bigtable) has a set of query requests. Each query is a set of boolean filters (clauses). We want to index a subset of fields (variables) to accelerate the queries. If we index a field, it satisfies positive/negative constraints. We want to maximize queries served.
\textbf{Model:} Max-SAT. Variables are index decisions, clauses are query filters. Randomized LP rounding finds the indexing strategy that maximizes query throughput.
\textbf{Why Randomized Rounding matters:} Real systems cannot wait for exact integer programming solvers. Randomized LP rounding runs instantly by solving a continuous linear program and sampling variable assignments, guaranteeing that query satisfaction is within a fraction of the absolute optimal configuration.

\section{Tight Example: Clause Size Bottleneck}
A tight example for the $(1-1/e)$ bound consists of $n$ independent single-literal clauses of the form $(x_i)$ for $i=1, \dots, n$, each with weight 1.
The LP relaxation will set $y_i^* = 1.0$ and $z_i^* = 1.0$, achieving objective $n$.
Randomized rounding sets each $x_i$ to True with probability $y_i^* = 1.0$, satisfying all clauses exactly.
Wait! What if we have a clause $(x_1 \lor x_2 \lor \dots \lor x_k)$? As $k$ grows, the factor $1 - (1 - 1/k)^k$ converges to $1 - 1/e \approx 0.632$. This shows the asymptotic bound is tight for large clauses under pure randomized rounding, but the coin-flip algorithm achieves $1 - 2^{-k} \approx 1$ for large clauses, which compensates for this bottleneck.

\section{Exercises}
\begin{exercise}
Prove that for any clause $C_j$ of size $k$, the probability that it is satisfied by the coin-flip assignment is $1 - 2^{-k}$.
\end{exercise}
\begin{exercise}
Show that if all clauses in a Max-SAT instance have size 1, the randomized rounding algorithm finds an exact optimal solution.
\end{exercise}

